\section{Rees charts, contact profiles and the marked module}\label{app:reesmarked}

Let \(A=\mathbb K[\lambda,\xi,\eta,\zeta]\) and
\(I=(\lambda,\xi\eta-\lambda\zeta)\).  The Rees presentation is
\[
A[Y_L,Y_Q]/((\xi\eta-\lambda\zeta)Y_L-\lambda Y_Q).
\]
On the \(Y_L\ne0\) chart, set \(r=Y_Q/Y_L\); the equation is
\(\xi\eta=\lambda(\zeta+r)\).  On \(Y_Q\ne0\), set \(s=Y_L/Y_Q\); the equation
is \(s(\xi\eta-\lambda\zeta)=\lambda\).  Above
\(\lambda=\xi=\eta=0\), both charts retain a positive-dimensional section,
which is why the blow-up does not separate the two components.

For an analytic arc, write
\(a=\ord\lambda,b=\ord\xi,c=\ord\eta\).  The orders of the union ideal and
the two component ideals are
\[
\ord I=\min(a,b+c),\quad
\ord(\lambda,\eta)=\min(a,c),\quad
\ord(\lambda,\xi)=\min(a,b).
\]
Adding the intersection order \(\min(a,b,c)\) gives \cref{eq:contact}.

The tagged ideal lives in \(A[u,x]\).  Its two-generator presentation
\[
A[u,x]^2\longrightarrow J,\qquad
(f,g)\longmapsto fLux^3+gQu^2x^4
\]
has kernel generated by \((Qux,-L)\).  Adding a formal \(R\)-basis vector
produces a three-generator marked module with the second kernel vector displayed
in \cref{sec:gf1}.  The section \(s_N\) keeps the resolvent-polynomial marking.
To obtain the amplitude section one must also multiply by the inverse
denominator series and by the depth-dependent phases \((-i)^k/k!\).

There is no claim that this marked module is a canonical physical state space.
It is an exact algebraic bookkeeping device for coefficient provenance.
